
This section devoted to the development of all the notions and results necessary to understand and prove the spectral theorem, stated below. 

\begin{theorem}[Spectral theorem]
For every self-adjoint operator $A\cl \mathcal{D}_A\to \mathcal{H}$ there is a unique projection-valued measure $\mathrm{P}_{\negmedspace A}\cl \sigma(\mathcal{O}_{\R})\to\mathcal{L}(\mathcal{H})$ such that
\begin{equation*}
A = \int_{\R}\! i_{\R} \, \d \mathrm{P}_{\negmedspace A} \equiv \int_{\R}\! \lambda \, \mathrm{P}_{\negmedspace A}(\d \lambda),
\end{equation*}
where $i_{\R}\cl \R\hookrightarrow \C$ is the inclusion of $\R$ into $\C$.
\end{theorem}

While useful in theory, existence results are often of limited use in practice since they usually only tell us that something exists, and not how to construct it. However, we should note here that the proof of the spectral theorem is, in fact, constructive in nature. Hence, given any self-adjoint operator $A$, we will be able to explicitly determine its associated projection-valued measure $\mathrm{P}_{\negmedspace A}$ along the following steps.
\begin{enumerate}[label=(\roman*)]
  \item For each $\psi\in\mathcal{H}$, construct the real-valued Borel measure $\mu^A_{\psi}\cl \sigma(\mathcal{O}_{\R})\to \R$ given by
  \begin{equation*}
  \mu^A_{\psi} ((-\infty,\lambda]):= \lim_{\delta\to 0^+} \lim_{\varepsilon\to 0^+} \int_{-\infty}^{\lambda+\delta}\!\d t  \Im\langle\psi | R_A(t+\mathrm{i}\varepsilon)\psi\rangle ,
  \end{equation*}
  where $R_A\cl\rho(A)\to \mathcal{L}(\mathcal{H})$ is the resolvent map of $A$. This is know as the \emph{Stieltjes inversion formula}. Note that while not every element in $\sigma(\mathcal{O}_{\R})$ is of the form $(-\infty,\lambda]$, such Borel measurable sets do generate the entire $\sigma(\mathcal{O}_{\R})$ via unions, intersections and set differences. Hence, the value of $\mu^A_{\psi}(\Omega)$ for $\Omega\in\sigma(\mathcal{O}_{\R})$ can be determined by applying the corresponding formulae for measures, namely $\sigma$-additivity, continuity from above and measure of set differences.
  \item For all $\psi,\varphi\in\mathcal{H}$, define the complex-valued Borel measure $\mu^A_{\psi,\varphi}\cl\sigma(\mathcal{O}_{\R})\to\C$ by
  \begin{equation*}
  \mu^A_{\psi,\varphi}(\Omega):=\tfrac{1}{4}(\mu^A_{\psi+\varphi}(\Omega)-\mu^A_{\psi-\varphi}(\Omega)+\mathrm{i}\mu^A_{\psi-\mathrm{i}\varphi}(\Omega)-\mathrm{i}\mu^A_{\psi+\mathrm{i}\varphi}(\Omega)).
  \end{equation*}
  \item Define the projection-valued measure $\mathrm{P}_{\negmedspace A}\cl \sigma(\mathcal{O}_{\R})\to\mathcal{L}(\mathcal{H})$ by requiring $\mathrm{P}_{\negmedspace A}(\Omega)$, for each $\Omega\in\sigma(\mathcal{O}_{\R})$, to be the unique map in $\mathcal{L}(\mathcal{H})$ satisfying
  \begin{equation*}
    \forall\, \psi,\varphi\in\mathcal{H}:\
    \langle\psi,\mathrm{P}_{\negmedspace A}(\Omega)\varphi\rangle
    = \int_{\R}\!\chi_{\Omega}\,\d\mu^A_{\psi,\varphi}.
    = \mu^{A}_{\psi,\varphi}(\Omega)
  \end{equation*}
\end{enumerate}

We will now make all the notions and constructions used herein precise. In fact, we will present the relevant definitions and results by taking the inverse route, starting with projection-valued measures and arriving at their associated self-adjoint operators, obtaining (and proving) what we will call the inverse spectral theorem.

\section{Projection-valued measures}

Projection-valued measures are, unsurprisingly, objects sharing
characteristics of both measures and projection operators.

\begin{definition}
  Let \(\Cal H\) be a Banach space and \(\Cal L(\Cal H)\) its algebra
  of bounded operators.
  A \emph{projection-valued measure}
  \index{projection-valued measure}
  on a measurable space \((M,\Sigma)\)
  (with values in \(\Cal H\)-projectors)
  is a function \(P\colon\Sigma\to\Cal L(\Cal H)\) satisfying
  the following properties:
  \begin{enumerate}[label=(\roman*)]
    \item
      $\forall \Omega \in \sigma(\mathcal{O}_{\R}), \
      P(\Omega)^* = P(\Omega) $
    \item
      $\forall \Omega \in \sigma(\mathcal{O}_{\R}), \
      P(\Omega)\circ P(\Omega) = P(\Omega) $
    \item \(P(M)=\id_{\Cal H}\).
    \item
      for any pairwise disjoint sequence
      $\{\Omega_n\}_{n\in\N}$ in $\Sigma$
      and for any \(\psi\in\Cal H\),
      \begin{equation*}
        \sum_{n=0}^{\infty}V(\Omega_n)\psi
        =
        V\Big(\bigcup_{\, n=0}^{\infty}\Omega_n\Big)\psi.
      \end{equation*}
  \end{enumerate}
\end{definition}

\begin{lemma}
  Let $P$ be an \(\Cal L(\Cal H)\)-valued
  measure on a measurable space \((M,\Sigma)\).
  Then, for any $\Omega,\Omega_1,\Omega_2\in\Sigma$,
  \begin{enumerate}[label=(\roman*)]
    \item
      \(P(\emptyset)=0\)
    \item
      $P(\Omega_1\cup\Omega_2)
      =P(\Omega_1)+P(\Omega_2)-P(\Omega_1\cap\Omega_2)$
    \item
      if $\Omega_1\subseteq\Omega_2$,
      then $\ran(P(\Omega_1))\subseteq \ran(P(\Omega_2))$.
    \item \(P(\Omega_1\cap\Omega_2)=P(\Omega_1)\circ P(\Omega_2)\).
  \end{enumerate}
\end{lemma}

Let \(P\colon\Sigma\to\Cal L(\Cal H)\) be a projection valued measure.
As with any operator, for each \(\Omega\), the information of
\(P(\Omega)\) is captured by the family of complex numbers
\begin{equation}
  \mu_{\psi,\phi}(\Omega) = \langle \psi,P(\Omega)\phi\rangle
\end{equation}
for each \(\psi,\phi\in\Cal H\).
By polarization, this is also captured by the family of positive real
numbers
\begin{equation}
  \mu_\psi(\Omega) = \langle \psi,P(\Omega)\psi\rangle
\end{equation}
for each \(\psi\in\Cal H\).
It turns out that each \(\mu_{\psi,\phi}\colon\Sigma\to\C\) is a
complex measure satisfying
\(\mu_{\psi,\phi}(M)=\langle\psi,\phi\rangle\), while each
\(\mu_\psi\colon\Sigma\to[0,\infty)\) is a
real positive measure satisfying \(\mu_\psi(M)=\|\psi\|^{2}\).

\section{Integration of bounded measurable functions}

We can integrate bounded functions \(f\colon M\to\C\) with respect to
a projection-valued measure \(P\), obtaining a bounded operator
\(\int_Mf\d P \in \Cal L(\Cal H)\).
Recall that a simple function \(s\colon M\to \C\) is a function
taking finitely many values, and that the space of simple measurable
functions \(\Cal S(M)\) is a dense linear subspace of the Banach space
\(\Cal B(M)\) of bounded measurable functions on \(M\).
This allows us to prove the following result.
\begin{proposition}
  The linear map
  \begin{align}
    \int_M\d P\colon \Cal S(M) &\to \Cal L(\Cal H) \\
    f &\mapsto \textstyle\int_M s\d P
  \end{align}
  defined,
  for a simple function \(s=\sum_{i=1}^{n}r_i\chi_{\Omega_i}\), as
  \begin{equation}
    \int_{M}^{} s \d P = \sum_{i=1}^{n} r_i P(\Omega_i),
  \end{equation}
  is a bounded linear map of norm \(1\).
  In particular, by the BLT theorem, it extends uniquely to a bounded
  linear map
  \begin{align}
    \int_M\d P\colon \Cal B(M) &\to \Cal L(\Cal H) \\
    f &\mapsto \textstyle\int_M f\d P
  \end{align}
  with norm \(1\).
\end{proposition}
\begin{proof}
   Using sudivisions to put \(s\) in a standard form, one readily
   checks that this is independent of the representation of \(s\) and
   that it is linear.
   To see that it is bounded, suppose that
   \(s=\sum_{i=1}^{n}r_i\chi_{\Omega_n}\) is in standard form (i.e.
   \(\Omega_i\cap\Omega_j=\emptyset\) whenever \(i\neq j\)).
   Since the \(P(\Omega_i)\) are self-adjoint, idempotent, pairwise
   orthogonal operators, we have
   \begin{align}
     \| (\textstyle\int_Ms\d P)\psi \|^{2}
     &= \|\textstyle\sum_{i=1}^{n}r_iP(\Omega_i)\psi \|^{2} \\
     &= \sum_{i,j=1}^{n}
       \langle r_iP(\Omega_i)\psi,r_jP(\Omega_j)\psi \rangle \\
     &= \sum_{i=1}^{n} |r_i|^{2}
       \langle \psi,P(\Omega_i)\psi \rangle \\
     &= \sum_{i=1}^{n} |r_i|^{2} \mu_\psi(\Omega_i) \\
     &= \int_{M} |s|^{2}\d\mu_\psi \\
     &\leq \|s\|_{\infty}^{2}\int_{M} \d\mu_\psi \\
     &\leq \|s\|_{\infty}^{2} \|\psi\|^{2}
   .\end{align}
   In fact, by choosing \(s=1\), we see that the bound is sharp.
   Hence, for each \(s\in\Cal S(M)\), the operator \(\int_Ms\d P\) is
   bounded with norm
   \begin{equation}
     \|\textstyle\int_Ms\d P\| = \|s\|_{\infty}
   .\end{equation}
   It follows that the operator \(s\mapsto \int_Ms\d P\) is bounded
   with norm \(1\).
\end{proof}

\iffalse
Sea \(\Cal H\) un espacio de Banach y
\(V\colon\Sigma\to\Cal H\) una medida vectorial sobre
un espacio medible \((M,\Sigma)\).
Para toda función simple medible
\(s=\Sigma_{i=1}^{n}r_i\chi_{\Omega_i}\), definamos
\begin{equation}
  \int_{M}s\d V = \sum_{i=1}^{n}r_iV(\Omega_i)
.\end{equation}
Esto define una función lineal \(\int_M\d V\colon \Cal S(M)\to\Cal
H\), donde \(\Cal S(M)\) es el espacio de las funciones simples
medibles. Pero no puedo probar que este operador sea acotado:
\begin{align}
  \| \textstyle\int_Ms\d V\|
  &= \|\textstyle\sum_{i=1}^{n}r_iV(\Omega_i)\| \\
  &\leq \sum_{i=1}^{n} |r_i|
    \| V(\Omega_i) \| \\
  &\leq \|s\|_\infty\sum_{i=1}^{n}
    \| V(\Omega_i) \| \\
  &\leq \dots ?
.\end{align}
\fi

\begin{proposition}
  \label{prop:propertiesint}
  The integral \(\int_M\d P\colon\Cal B(M)\to\Cal L(\Cal H)\) is a
  linear map satisfiyng:
  \begin{enumerate}
    \item
      For any bounded \(f\colon M\to\C\) and any \(\psi\in\Cal 
      H\),
      \begin{equation}
        \label{eq:fmupsi}
        \Big\langle \psi,\Big(\int_Mf\d P\Big) \psi\Big\rangle
        =
        \int_M f\d\mu_\psi
      .\end{equation}
    \item 
      For any \(\Omega\in\Sigma\),
      \begin{equation}
        \int_M \chi_\Omega \d P = P(\Omega)
      .\end{equation}
    \item 
      For any bounded \(f,g\colon M\to\C\),
      \begin{equation}
        \int_M fg \d P = 
        \Big(\int_M f \d P\Big)
        \Big(\int_M g \d P \Big) 
      \end{equation}
    \item 
      For any bounded \(f\colon M\to\C\)
      \begin{equation}
        \int_M \bar f \d P = 
        \Big(\int_M f \d P\Big)^{*}.
      \end{equation}
    \item 
      and
      \begin{equation}
        \Big\|\int_Mf\d P\Big\| \leq \|f\|_{\infty}
      \end{equation}
  \end{enumerate}
\end{proposition}
\begin{proof}
  We only show the first one.
  For a simple function \(s=\sum_{i=1}^{n}r_i\chi_{\Omega_i}\), we
  have
  \begin{align}
    \Big\langle \psi,\Big(\int_Ms\d P\Big) \psi\Big\rangle
    &= \Big\langle \psi,\sum_{i=1}^{n}r_iP(\Omega_i)\psi\Big\rangle \\
    &= \sum_{i=1}^{n}r_i\langle \psi,P(\Omega_i)\psi\rangle \\
    &= \sum_{i=1}^{n}r_i \mu_\psi(\Omega_i) \\
    &= \int_Ms\d P
  .\end{align}
  Now let \(f\) be a bounded function and \(s_n\) a sequence of simple
  functions converging to \(f\) in \(\Cal B(M)\). By bounded
  convergence and continuity of the inner product, we get
  \begin{align}
    \Big\langle \psi,\Big(\int_Mf\d P\Big) \psi\Big\rangle
    &= \Big\langle \psi,
      \lim_{n\to\infty}
      \Big(\int_Ms_i\d P\Big) \psi\Big\rangle \\
    &= \lim_{n\to\infty}
      \Big\langle \psi,\Big(\int_Ms_i\d P\Big) \psi\Big\rangle \\
    &= \lim_{n\to\infty} \int_Ms_i\d P \\
    &= \int_Mf\d P
  .\end{align}
\end{proof}

\section{Integration of measurable functions}

So far we have obtained an integral that takes each bounded function
\(f\colon M\to\C\) to a bounded operator \(\int_Mf\d P\in\Cal L(\Cal
H)\).
However, for the spectral theorem, we will need to integrate arbitrary
measurable functions \(f\colon M\to\Cal C\). We can do it, except now
the result \(\int_Mf\d P\) will be a (possibly unbounded) densely
defined operator \(\Cal D_f\to\Cal H\).

Suppose we want this extended definition to preserve the first four
properties in \eqref{eq:fmupsi}. Then we have
\begin{align}
  \Big\|\Big(\int_Mf\d P\Big) \psi\Big\|^{2}
  &= \Big\langle\Big(\int_Mf\d P\Big) \psi,
    \Big(\int_Mf\d P\Big) \psi\Big\rangle \\
  &= \Big\langle \psi,
    \Big(\int_M\bar f f\d P\Big) \psi\Big\rangle \\
  &= \int_M |f|^{2}\d\mu_\psi
.\end{align}
If this has any hope of making sense, then it means that the right
hand side must be finite.
Hence, given an arbitrary measurable \(f\colon M\to\C\),
we restrict the domain of \(\int_Mf\d P\) to be those \(\psi\)
such that \(f\) is square \(\mu_\psi\)-integrable:
\begin{equation}
  \Cal D_{\int_Mf\d P}
  =
  \left\{
    \psi\in\Cal H
    :
    f \in L^{2}(\mu_\psi)
  \right\}
.\end{equation}
It takes a bit of work to prove that \(\Cal D_{\int_Mf\d P}\) is
indeed a dense subspace of \(\Cal H\), which we won't do here.

Now we define the operator \(\int_Mf\d P\).
Observe that, given \(\phi,\psi\in\Cal D_{\int_Mf\d P}\) the
function \(f\) is also square-integrable with respect to the measure
\begin{equation}
  \mu_{\phi,\psi}
  = \frac{1}{4}
  (\mu_{\phi+\psi}-\mu_{\phi-\psi}+\ic\mu_{\phi-\ic\psi}-\ic\mu_{\phi+\ic\psi})
\end{equation}
since it is integrable with respect to each summand.
Fix \(\psi\). Then for each \(\phi\), the conjugate-linear map
\begin{align}
  \label{eq:integrationpsi}
  \Cal D_{\int_Mf\d P} & \to \C \\
  \phi&\mapsto\textstyle\int_{M}f\d\mu_{\phi,\psi}
\end{align}
is bounded: indeed, fix a sequence of bounded maps \(f_n\) converging to
\(f\) pointwise.
Then for every \(\phi\in\Cal D_{\int_Mf\d P}\), \(f\) is also in
\(L^{2}(\mu_{\phi,\psi})\), so we get
\begin{align}
  \Big |\int_{M}f\d\mu_{\phi,\psi} \Big |
  &= \lim_{n\to\infty}\Big |\int_{M}f_n\d\mu_{\phi,\psi} \Big | \\
  &= \lim_{n\to\infty}
  \Big |\Big\langle
    \phi ,\Big(\int_{M}f_n\d P\Big)\psi
  \Big\rangle\Big | \\
  &\leq \lim_{n\to\infty}
    \|\phi\|\Big\|\Big(\int_{M}f_n\d P\Big)\psi\Big\| \\
  &= \lim_{n\to\infty}
    \|\phi\|\Big\|\int_{M}f_n\d\mu_\psi\Big\| \\
  &= \|\phi\|\Big\|\int_{M}f\d\mu_\psi\Big\| \\
  &= \|\phi\| \cdot \|f\|_{L^{2}(\mu_\psi)}.
\end{align}
By the BLT theorem, the map \eqref{eq:integrationpsi} 
extends to a bounded conjugate-linear map
\(\Cal H\to\C\) which, by the (conjugate) Riesz representation theorem,
is of the form \(\langle-,\eta\rangle\) for a unique \(\eta\), which
we call \((\int_Mf\d P)\psi\).
In particular, for \(\phi\) in the domain, we have
\begin{equation}
  \Big\langle \phi, \Big(\int_Mf\d P\Big)\psi\Big\rangle
  =
  \int_{M}f\d\mu_{\phi,\psi}
\end{equation}
Moreover, the assignment \(\psi\mapsto
(\int_Mf\d P)\psi\) is linear, so it defines our desired operator
\begin{equation}
  \int_Mf\d P \colon \Cal D_{\int_Mf\d P} \to \Cal H.
\end{equation}

Since the set of densely defined operators is
not a vector space, it makes no sense to ask if the integration
operator \(f\mapsto \int_Mf\d P\) is linear.
Nevertheless, we have we still have
\(\int_{M}^{}\alpha f\d P = \alpha \int_{M}^{}f\d P\), but the
integral is compatible with sums only to a certain degree:
\begin{equation}
  \int_{M}f\d P
  +
  \int_{M}g\d P
  \subseteq 
  \int_{M}(f+g)\d P
.\end{equation}
For unbounded functions \(f,g\), property 3 in
\cref{prop:propertiesint} must be replaced by
\begin{equation}
  \int_{M}f\d P
  \circ
  \int_{M}g\d P
  \subseteq 
  \int_{M}fg\d P
.\end{equation}
Importantly, property 4 is still valid for unbounded functions.
A consequence is that, if \(f\colon M\to\R\) is a real,
measurable function, then its integral \(\int_Mf\d P\) is a
\emph{self-adjoint} operator.
An even more special case of this observation is the following
corollary.
\begin{corollary}
  [Inverse spectral theorem]
  \label{cor:inverse-spectral-theorem}
  Given a projection-valued measure \(P\colon\Sigma_\R\to\Cal L(\Cal H)\)
  on the Borel sigma algebra of \(\R\), the integral
  \begin{equation}
    A_P = \int_{\R}i_{\R}\d P
  \end{equation}
  is a self-adjoint operator, where \(i_\R\) is the natural inclusion
  of \(\R\) into \(\C\).
\end{corollary}



